In this section, we briefly connect the IJ covariance to some related
literature on local sensitivity, outlier detection, and
the von Mises calculus.  The connections are illuminated by the following
``weighted posterior expectation''
%
\begin{align}
    %
    \expect{\postw}{\g(\theta)} =
        \frac{
        \int_{\thetadom} g(\theta)
            \exp\left(N \meann \w_n \ell(\x_n \vert \theta) \right) \prior(\theta)
                            d\theta}
             {\int_{\thetadom}
             \exp\left(N \meann \w_n \ell(\x_n \vert \theta) \right) \prior(\theta)
             d\theta},
             \eqlabel{weighted_post}
    %
\end{align}
%
where $\wvec = (\w_1, \ldots, \w_N) \in \rdom{N}$ is a length--$N$ vector
of scalar--valued weights, one for each datapoint.  Comparing
\eqref{g_post} with \eqref{weighted_post}, we see that
$\expect{\postw}{\g(\theta)} = \expect{\post}{\g(\theta)}$ when
$\wvec = \onevec$, where $\onevec$ is the length--$N$ vector containing
all ones.  When we can apply the dominated convergence theorem to
exchange posterior integration and differentiation
(see, e.g., Theorem 1 of \citet{giordano:2018:covariances}), it follows that
%
\begin{align}
    \fracat{\partial \expect{\postw}{\g(\theta)}}{\partial \w_n}{\wvec = \onevec} =
    \cov{\post}{\g(\theta), \ell(\x_n \vert \theta)} = \frac{1}{N} \infl_n.
    \eqlabel{infl_deriv}
\end{align}
%
Our ``influence scores'' are thus the values of the
``empirical influence function'' of the statistical functional $\expect{\post}{\g(\theta)}$,
and our $\gcovij$ the classical IJ covariance estimator
(see, e.g., \citet{mises:1947:asymptotic,reeds:1976:thesis,efron:1982:jackknife,hampel:2011:robustbook}, or
\citet{serfling:1980:approximation} for a textbook treatment).  Equivalently,
the empirical influence function can be understood as an instance of
``local sensitivity analysis'' to perturbations of data weights
\citep{gustafson:1996:localposterior,giordano:2018:covariances}, which
motivates application of the empirical influence function
to a wider variety of problems than variance estimation.
The related literature is too large to survey here, but
some examples are cross--validation
\citep{koh:2017:blackboxinfluence,giordano:2019:swiss}, identification
of outliers \citep{pregibon:1981:logistic,mcculloch:1989:localbayesianmodelinfluence,belsley:1980:regression},
identification of influential subsets \citep{giordano:2026:amip1,giordano:2026:amip2},
and interpretation of Bayesian posteriors
\citep{thomas:2018:reconcilingcurvatureandis,iba:2023:wkernel}.  The
present work rigorously connects Bayesian posterior expectations to this
large literature, and we anticipate applications beyond variance estimation.

However, the interpretation of $\infl_n$ as a local approximation, together
with the form of \eqref{weighted_post}, provides a cautionary warning
in the Bayesian context.  Though our experiments of \secref{experiments}
motivate the IJ covariance estimator as a way to avoid pathologies associated
with MAP estimation, the most common use of MCMC is to integrate out
high--dimensional posterior parameters that do not obey a BCLT.  The following
simple argument suggests that the IJ covariance may not work in such a case,
even for parameters that do obey a BCLT marginally.

Suppose, for example, that $\theta$ is the concatenation of two sets of
parameters: a finite--dimensional set of ``global'' parameters $\gamma$, and a
high--dimensional set of ``local'' parameters, $\lambda$.  A common such example
is where $\gamma$ are ``fixed effects'' and $\lambda$ ``random effects'' in a
panel data setting, and the full parameter set is the concatenation $\theta =
(\gamma, \lambda)$ (see, e.g.~\citet[Chapter 7]{davidson:2004:econometric},
though in a Bayesian context the distinction between fixed and random effects is
not always meaningful \citep[Chapter 11]{gelman:2006:arm}).  Suppose that
$\cov{\post}{\gamma} = \ordlogp{N^{-1}}$ but $\cov{\post}{\lambda}$ remains
bounded away from zero.  Take $\g(\theta) = \gamma$. Differentiating
\eqref{infl_deriv} again with respect to $\w_n$ and applying the tower property
and delta method gives:
%
\begin{align}
    \fracat{\partial^2 \expect{\postw}{\gamma}}{\partial \w_n^2}{\wvec = \onevec} =
    % \expect{\post}{
    %     \gammabar \ellbar(\x_n \vert \gamma, \lambda)^2
    %     }
    % =
    \cov{\p(\gamma \vert \xvec)}{
        \gamma,
        \cov{\p(\lambda \vert \gamma, \xvec)}{\ell_n(\gamma, \lambda)}
    } + \ordlogp{N^{-2}}
    \eqlabel{resid_term}
\end{align}
%
% where for compactness we have written $\gammabar = \gamma - \expect{\post}{\gamma}$
% and $\ellbar(\x_n \vert \theta) = \ell(\x_n \vert \theta) - \expect{\post}{\ell(\x_n \vert \theta)}$
%
where the $\ordlogp{N^{-2}}$ term assumes that third--order posterior cumulants
of $\p(\gamma \vert \xvec)$ are $\ordlogp{N^{-2}}$, as they would typically be
for a posterior which concentrates at root $N^{-1/2}$. By the tower property,
the first derivative is also a posterior covariance with respect to  $\p(\gamma
\vert \xvec)$:
%
\begin{align}
    \fracat{\partial \expect{\postw}{\gamma}}{\partial \w_n}{\wvec = \onevec} =
    \cov{\p(\gamma \vert \xvec)}{
        \gamma,
        \expect{\p(\lambda \vert \xvec, \gamma)}{\ell(\x_n \vert \gamma, \lambda)}
    }.
    \eqlabel{global_local_first_order}
\end{align}
%
See \appref{weights_higher_order} for details on the derivation of
\eqref{resid_term,global_local_first_order}.


If the function $\gamma \mapsto \cov{\p(\lambda \vert \gamma,
\xvec)}{\ell_n(\gamma, \lambda)}$ does not vanish as $N \rightarrow \infty$, it
follows from \eqref{resid_term,global_local_first_order} that both the first and
second derivatives are posterior covariances of non--vanishing functions of
$\gamma$, and so of the same order as $N \rightarrow \infty$.  Since the second
derivative is of the same order as the first derivative, the empirical influence
function's linear approximation to provide a poor approximation even to the
effect of leaving out a single datapoint.

Note that this counterexample
does not contradict our results from \thmref{ij_consistent}: if a BCLT applies
to $\lambda$ as well as to $\gamma$, then $\cov{\p(\lambda \vert \xvec,
\gamma)}{\ell(\x_n \vert \gamma, \lambda)} = \ordlogp{N^{-1}}$ inside the
covariance of \eqref{resid_term}, and the second derivative is
$\ordlogp{N^{-2}}$, as suggested by \thmref{bayes_clt_main}.
%
This intuition is rigorously developed and supported using von Mises expansions
in \citet{giordano:2024:oldbayesij}, which at the time of writing is still in
progress.
In light of this example, and of the analysis of
\citet{giordano:2024:oldbayesij} for the high--dimensional regime, we recommend
caution when employing the IJ covariance estimator --- or any method depending
on the empirical influence function ---  in settings where a BCLT may not apply
to the full posterior $\post$.


